\documentclass[12pt,a4paper]{article}
\usepackage{ugmath}
\usepackage{float}
\usepackage{placeins}
\usepackage [spanish]{babel}
\usepackage[labelfont=bf,labelsep=space]{caption}
\newcommand{\paren}[1]{\left( #1 \right)}
\newcommand{\alumno}{Ricardo León Martínez}
\newcommand{\materia}{Calculo Diferencial e Integral II}
\newcommand{\profesor}{Fernando Nuñez Medina}
\newcommand{\tarea}{Tarea 10}
\newcommand{\fecha}{2/5/2026}

\begin{document}
    \begin{center}
    {\large\textbf{UNIVERSIDAD DE GUANAJUATO}}\\[0.3cm]
    {\normalsize\textbf{DIVISIÓN DE CIENCIAS NATURALES Y EXACTAS}}\\
    {\normalsize\textbf{CAMPUS GUANAJUATO}}\\[1cm]

    {\Large\textbf{\tarea\ (\materia)}}\\[1cm]
    \end{center}

    \textbf{Nombre:} \alumno \hfill 
    \textbf{Fecha:} \fecha \hfill 
    \textbf{Calificación:} \rule{3cm}{0.4pt} \\[0.3cm]

    \begin{mdframed}[style=mdpinkbox,frametitle={Ejercicio 1}]
        \textbf{(Identidades del ángulo medio)} Prueba lo siguiente:
        \begin{enumerate}
            \item[(a)] $\sin(\frac{x}{2})=\pm\sqrt{\frac{1-\cos(x)}{2}}$.
            \item[(b)] $\cos(\frac{x}{2})=\pm\sqrt{\frac{1+\cos(x)}{2}}$.
            \item[(c)] $\tan(\frac{x}{2})=\pm\sqrt{\frac{1-\cos(x)}{1+\cos(x)}}$.  
        \end{enumerate}
    \end{mdframed}

    \begin{proof}
        \begin{enumerate}
            \item[(a)] $\sin(\frac{x}{2})=\pm\sqrt{\frac{1-\cos(x)}{2}}$.\\
            Por el corolario 3 sabemos que
            \begin{equation*}
                \sin^{2}(\theta)=\frac{1-\cos(2\theta)}{2},\quad\forall\theta\in\R.
            \end{equation*}
            En particular si $\theta=\frac{x}{2}$, entonces
            \begin{equation*}
                \sin^{2}\paren{\frac{x}{2}}=\frac{1-\cos(x)}{2}.
            \end{equation*}
            Aplicando la raiz cuadrada a ambos lados, obtenemos
            \begin{equation*}
                \abs{\sin\paren{\frac{x}{2}}}=\sqrt{\frac{1-\cos(x)}{2}}.
            \end{equation*}
            Por la definición de valor absoluto concluimos que
            \begin{equation*}
                \sin(\frac{x}{2})=\pm\sqrt{\frac{1-\cos(x)}{2}}.
            \end{equation*}
            \item[(b)] $\cos(\frac{x}{2})=\pm\sqrt{\frac{1+\cos(x)}{2}}$.\\
            Por el corolario 3 sabemos que
            \begin{equation*}
                \cos^{2}(\theta)=\frac{1+\cos(2\theta)}{2},\quad\forall\theta\in\R.
            \end{equation*}
            En particular si $\theta=\frac{x}{2}$, entonces
            \begin{equation*}
                \cos^{2}\paren{\frac{x}{2}}=\frac{1+\cos(x)}{2}.
            \end{equation*}
            Aplicando la raiz cuadrada a ambos lados, obtenemos
            \begin{equation*}
                \abs{\cos\paren{\frac{x}{2}}}=\sqrt{\frac{1+\cos(x)}{2}}.
            \end{equation*}
            Por la definición de valor absoluto concluimos que
            \begin{equation*}
                \cos(\frac{x}{2})=\pm\sqrt{\frac{1+\cos(x)}{2}}.
            \end{equation*}
            \item[(c)] $\tan(\frac{x}{2})=\pm\sqrt{\frac{1-\cos(x)}{1+\cos(x)}}$.\\
            Por el corolario 3 sabemos que
            \begin{equation*}
                \tan^{2}(x)=\frac{1-cos(2\theta)}{1+\cos(2\theta)},\quad\forall\theta\in\R.
            \end{equation*}
            En particular si $\theta=\frac{x}{2}$, entonces
            \begin{equation*}
                \tan^{2}\paren{\frac{x}{2}}=\frac{1-\cos(x)}{1+\cos(x)}.
            \end{equation*}
            Aplicando la raiz cuadrada a ambos lados, obtenemos
            \begin{equation*}
                \abs{\tan\paren{\frac{x}{2}}}=\sqrt{\frac{1-\cos(x)}{1+\cos(x)}}.
            \end{equation*}
            Por la definición de valor absoluto concluimos que
            \begin{equation*}
                \tan\paren{\frac{x}{2}}=\pm\sqrt{\frac{1-\cos(x)}{1+\cos(x)}}.
            \end{equation*}
        \end{enumerate}
    \end{proof}

    \begin{mdframed}[style=mdpinkbox,frametitle={Ejercicio 2}]
        Considera una sucesión $(a_{n})$. Prueba que si $a_{n}\to\pm\infty$ y $(a_{n_{k}})$
        es una subsucesión de $(a_{n})$, entonces $(a_{n_{k}})\to\pm\infty$.
    \end{mdframed}

    \begin{proof}
        Por definición de divergencia a $+\infty$
        \begin{equation*}
            \exists N\in\N\text{ tal que }\forall n\geq N,\quad a_{n}\gt M.
        \end{equation*}
        Como $n_{k}\to\infty$ cuando $k\to\infty$, podemos elegir $K\in\N$ tal que
        \begin{equation*}
            \forall k\geq K,\quad n_{k}\geq N.
        \end{equation*}
        Entonces para todo $k\geq K$ tenemos
        \begin{equation*}
            a_{n_{k}}\gt M.
        \end{equation*}
        Así,
        \begin{equation*}
            \forall M\gt0,\exists K\in\N\text{ tal que }\forall k\geq K, a_{n_{k}}\gt M,
        \end{equation*}
        que es precisamente $(a_{n_{k}})\to+\infty$. El caso para $a_{n}\to-\infty$ es analogo.
    \end{proof}

    \begin{mdframed}[style=mdpinkbox,frametitle={Ejercicio 3}]
        \textbf{(Agrupamiento de una serie)} Dada una serie $\sum_{n=1}^{\infty}a_{n}$ y una sucesión
        creciente de números naturales $n_{1}\lt n_{2}\lt n_{3}\lt\cdots$, definamos
        \begin{align*}
            A_{1} & =a_{1}+\cdots a_{n_{1}},\\
            A_{2} & =a_{n_{1}+1}+\cdots+a_{n_{2}},\\
            A_{3} & =a_{n_{2}+1}+\cdots+a_{n_{3}},\\
            & \vdots\\
            A_{k} & =a_{n_{k-1}+1}+\cdots+a_{n_{k}}.
        \end{align*}
        A la serie $\sum_{k=1}^{\infty}A_{k}$ le llamaremos un agrupamiento de la serie $\sum_{n=1}^{\infty}a_{n}$.
        Prueba que si la serie $\sum_{n=1}^{\infty}a_{n}$ converge a $L$, entonces todos sus agrupamientos también
        convergen a $L$.
    \end{mdframed}

    \begin{mdframed}[style=mdpinkbox,frametitle={Ejercicio 4}]
        Determina la convergencia de las series siguientes:
        \begin{multicols}{2}
            \begin{enumerate}
                \item[(a)]  $\sum_{n=1}^{\infty}$$\frac{n}{n+1}$.
                \item[(b)]  $\sum_{n=1}^{\infty}\frac{2^{n}+3^{n}}{4^{n}}.$
                \item[(c)]  $\sum_{n=1}^{\infty}\frac{5n}{n^{2}+1}.$
                \item[(d)]  $\sum_{n=2}^{\infty}\frac{\text{cos}\left(n^{3}+9\right)}{n^{3}+9}.$
            \end{enumerate}
        \end{multicols}
    \end{mdframed}

    \begin{solucion}
        \begin{enumerate}
            \item[(a)]  $\sum_{n=1}^{\infty}$$\frac{n}{n+1}$.
            \begin{equation*}
                \frac{n}{n+1}=1-\frac{1}{n+1}\to1\neq0\text{ cuando }n\to\infty.
            \end{equation*}
            Como el término general no tiende a 0, la serie diverge.
            \item[(b)]  $\sum_{n=1}^{\infty}\frac{2^{n}+3^{n}}{4^{n}}.$
            \begin{equation*}
                \frac{2^{n}+3^{n}}{4^{n}}=\paren{\frac{1}{2}}^{n}+\paren{\frac{3}{4}}^{n}
            \end{equation*}
            Ambas son geométricas de razón $r_{1}=\frac{1}{2}\lt1$ y $r_{2}=\frac{3}{4}\lt1$,
            luego convergen. Suma
            \begin{align*}
                \sum_{n=1}^{\infty}\paren{\frac{1}{2}}^{n}=\frac{\frac{1}{2}}{1-\frac{1}{2}}=1,\\
                \sum_{n=1}^{\infty}\paren{\frac{3}{4}}^{n}=\frac{\frac{3}{4}}{1-\frac{3}{4}}=3.
            \end{align*}
            Por lo tanto, la serie converge y su suma es $1+3=4$.
        \end{enumerate}
    \end{solucion}

    \begin{mdframed}[style=mdpinkbox,frametitle={Ejercicio 5}]
        \textbf{(Expansión decimal)}
        \begin{enumerate}
            \item[(a)] Prueba que si $(a_{n})$ es una sucesión de números enteros con $0\leq a_{n}\leq9$,
            entonces la serie
            \begin{equation*}
                \sum_{n=1}^{\infty}\frac{a_{n}}{10^{n}}
            \end{equation*}
            converge a un número del intervalo [0,1].
            \item[(b)] Prueba que si $x\in[0,1]$, entonces existe una sucesión de números enteros $(a_{n})$
            con $0\leq a_{n}\leq9$ tal que
            \begin{equation*}
                x=\sum_{n=1}^{\infty}\frac{a_{n}}{10^{n}}.
            \end{equation*}
            \item[(c)] Prueba que si $x\geq0$, entonces existen un entero no positivo $N$ y una sucesión de números
            enteros $(a_{n})_{n=N}^{\infty}$ con $0\leq a_{n}\leq9$ tales que
            \begin{equation*}
                x=\sum_{n=N}^{\infty}\frac{a_{n}}{10^{n}}.
            \end{equation*}
        \end{enumerate}
    \end{mdframed}

    \begin{proof}
        \begin{enumerate}
            \item[(a)]
            Sea $(a_{n})$ una sucesión con $0\leq a_{n}\leq 9$. Consideremos la serie
            \begin{equation*}
                \sum_{n=1}^{\infty}\frac{a_{n}}{10^{n}}.
            \end{equation*}
            Para todo $n$,
            \begin{equation*}
                0\leq\frac{a_{n}}{10^{n}}\leq\frac{9}{10^{n}}.
            \end{equation*}
            La serie geométrica
            \begin{equation*}
                \sum_{n=1}^{\infty}\frac{9}{10^{n}}
            \end{equation*}
            converge, pues la razón $r=\frac{1}{10}\lt1$, y la suma es
            \begin{equation*}
                \sum_{n=1}^{\infty}\frac{9}{10^{n}}=9\cdot\frac{\frac{1}{10}}{1-\frac{1}{10}}=1.
            \end{equation*}
            Por el criterio de comparación, la serie dada converge. Ademas,
            \begin{equation*}
                0\leq\sum_{n=1}^{\infty}\frac{a_{n}}{10^{n}}\leq\sum_{n=1}^{\infty}\frac{9}{10^{n}}=1.
            \end{equation*}
            Por lo tanto, la suma pertenece a $[0,1]$.
            \item[(b)]
            Sea $x\in[0,1]$. Definimos una sucesión $(a_{n})$ y residuos $(r_{n})$ como sigue: Sea $r_{0}=x$. Para $n\geq1$,
            \begin{equation*}
                a_{n}:=10r_{n-1},\quad r_{n}:=10r_{n-1}-a_{n}.
            \end{equation*}
            Entonces $0\leq a_{n}\leq9$, pues $0\leq r_{n-1}\leq1$. $0\leq r_{n}\leq1$. Iterando
            \begin{equation*}
                x=\frac{a_{1}}{10}+\frac{r_{1}}{10}=\frac{a_{1}}{10}+\frac{a_{2}}{10^{2}}+\frac{r_{2}}{10^{2}}=\cdots
            \end{equation*}
            y en general,
            \begin{equation*}
                x=\sum_{k=1}^{n}\frac{a_{k}}{10^{k}}+\frac{r_{n}}{10^{n}}.
            \end{equation*}
            Como $0\leq r_{n}\lt1$, se tiene
            \begin{equation*}
                0\leq\frac{r_{n}}{10^{n}}\leq\frac{1}{10^{n}}\to0.
            \end{equation*}
            Por lo tanto,
            \begin{equation*}
                x=\sum_{n=1}^{\infty}\frac{a_{n}}{10^{n}}.
            \end{equation*}
            \item[(c)]
            Sea $x\geq0$. Existe un entero $m\geq0$ tal que
            \begin{equation*}
                10^{m}\leq x\leq10^{m+1}.
            \end{equation*}
            Definimos
            \begin{equation*}
                y:=\frac{x}{10^{m+1}}\in[0,1).
            \end{equation*}
            Por (b), existe $(a_{n})_{n=1}^{\infty}$ con $0\leq a_{n}\leq9$ tal que
            \begin{equation*}
                y=\sum_{n=1}^{\infty}\frac{a_{n}}{10^{n}}.
            \end{equation*}
            Entonces,
            \begin{equation*}
                x=\sum_{n=1}^{\infty}\frac{a_{n}}{10^{n-(m+1)}}.
            \end{equation*}
            Definiendo $N:=-m$ y reindexando,
            \begin{equation*}
                x=\sum_{n=N}^{\infty}\frac{a_{n}}{10^{n}},
            \end{equation*}
            donde $0\leq a_{n}\leq9$.
        \end{enumerate}
    \end{proof}
\end{document}